\chapter{Literature Review}
\label{ch:literature}
This chapter summarises the principles of SAR image formation and the effects of terrain on SAR images that the
simulator has to reproduce, and reviews the existing approaches to SAR simulation and DEM-based geocoding.

\section{Range Resolution and Pulse Compression}
A SAR transmits a linear frequency-modulated pulse (chirp) $s(t)=\exp(j\pi K_r t^2)$ for $|t|<T_p/2$, of bandwidth
$B=K_rT_p$. A scatterer at slant range $R$ returns a copy of the pulse delayed by $2R/c$ and phase-shifted by
$-4\pi R/\lambda$. Correlating the echo with the transmitted replica (matched filtering) compresses the long pulse into a
sinc-shaped peak whose $-3$\,dB width defines the slant-range resolution
\begin{equation}
\rho_r = 0.886\,\frac{c}{2B}, \qquad \rho_g = \frac{\rho_r}{\sin\theta},
\end{equation}
where $\rho_g$ is the corresponding ground-range resolution at incidence angle $\theta$ \cite{cumming2005,curlander1991}.
The resolution therefore depends on the bandwidth and not on the pulse length, which can be chosen long enough to
deliver the required energy.

\section{Azimuth Resolution and the Synthetic Aperture}
Along the flight track (azimuth) a single pulse cannot separate targets, since the real antenna beam is several
kilometres wide on the ground. As the platform moves with speed $v$, the range to a target follows the hyperbola
$R(\eta)=\sqrt{R_0^2+v^2\eta^2}$, which produces a phase history $\exp(-j4\pi R(\eta)/\lambda)$ that is a chirp in
Doppler with a rate of $2v^2/(\lambda R_0)$. Matched filtering in azimuth over the synthetic aperture
$L_{sa}=\lambda R_0/L$, where $L$ is the antenna length, yields an azimuth resolution of about $L/2$, independent of
range \cite{moreira2013,curlander1991}. Several focusing algorithms exist; time-domain back-projection, in which the
range-compressed samples at the exact range of every image point are phase-corrected and summed over all pulses, is
the most direct and is exact for arbitrary scene geometry \cite{cumming2005}.

\section{Speckle}
A natural surface contains many elementary scatterers per resolution cell. Their coherent sum has a random phase, so
that the single-look intensity of a homogeneous surface is exponentially distributed with a standard deviation equal to
its mean \cite{goodman1976}. Averaging $L$ independent looks (multilooking) produces a Gamma distribution with a
relative variance of $1/L$, at the cost of spatial resolution \cite{oliver1998}. Speckle is therefore a property of the
measurement rather than of the surface, and comparisons between simulated and real images are made on multilooked
data.

\section{Geometric Distortions in Mountainous Terrain}
Because a SAR orders echoes by slant range, terrain height $h$ shifts a point towards the sensor by $h\cot\theta$ in
ground range, and the slope of the terrain controls how much ground falls into each range cell
\cite{leberl1990}. Three regimes follow (Fig.~\ref{fig:intro}). Slopes facing the radar with a slope smaller than the
look angle are compressed into fewer range cells and appear bright (foreshortening). If the slope exceeds the look
angle, the top of the slope is closer to the radar than its foot, so that the slope is imaged reversed and mixed with
the terrain in front of it (layover). Slopes facing away from the radar with a slope larger than the complement of the
look angle, together with all terrain hidden behind them, return no echo (radar shadow). Kropatsch and Strobl
\cite{kropatsch1990} showed how layover and shadow maps can be derived from a DEM, and Gelautz et al.
\cite{gelautz1998} analysed simulated and real SAR images of Alpine terrain, where these effects are pervasive.

\section{Radiometry and Backscatter Models}
The radar cross-section per unit ground area, $\sigma^0$, depends on the surface material, its roughness and moisture,
and on the local incidence angle $\theta_{loc}$ between the radar beam and the surface normal
\cite{ulaby1982,ulaby1989}. Simulators commonly reduce this dependence to a simple law of $\theta_{loc}$, such as the
cosine law, Lambert's law ($\cos^2\theta_{loc}$) or Muhleman's law \cite{muhleman1964}, whose validity for real
terrain remains an empirical question. In addition, the area of terrain that maps into one image pixel varies with the
topography. ESA calibrates GRD products to $\sigma^0$ referred to the ellipsoid pixel area, so that the brightness of a
pixel contains both effects; separating them is the purpose of radiometric slope correction and terrain flattening
\cite{ulander1996,small2011,frey2013}.

\section{SAR Simulators and DEM-Based Geocoding}
SAR simulators span a wide range of complexity. Raw-signal simulators generate the echo of every pulse and are used to
test focusing algorithms and system designs \cite{franceschetti1992}. Ray tracers model multiple bounces in urban
scenes and are used to interpret high-resolution images of buildings \cite{balz2009,auer2010}. Image-domain
simulators, finally, project each DEM facet directly into the image and form the basis of operational geocoding and
terrain correction \cite{schreier1993,meier1993,small2011,low2007}; they are also used for DEM-assisted co-registration
in interferometry \cite{nitti2011,bamler1998}. The absolute geolocation accuracy of Sentinel-1 products, which bounds
how well any simulation can be aligned with them without co-registration, has been verified with corner reflectors and
reaches a few metres when atmospheric delay and solid-earth tides are corrected \cite{schubert2015,gisinger2021}.

\section{Research Gap}
Existing studies either use DEM-based simulation as an internal step of geocoding and terrain-correction chains, or
analyse simulated images qualitatively. A transparent simulator that is verified from its physical principles upwards,
driven by the exact annotation of current Sentinel-1 products, and assessed quantitatively against real images for
several backscatter laws, viewing geometries, processing choices and land-cover classes is not readily available. This
work addresses this gap.
